% =====================================================================
% PROJECT ONE-PAGER TEMPLATE
% Matches the visual style of latex-resumes/resume.tex
%
% HOW TO REUSE PER PROJECT:
%   1. Copy this file and rename it (e.g. project-routerunner.tex)
%   2. Fill in every [BRACKETED PLACEHOLDER] below
%   3. Keep bullets short -- this is a single page, no more than
%      ~4 Approach bullets and 1-3 Result bullets will comfortably fit
%   4. Delete the \href / icon lines you don't need (e.g. no demo link)
% =====================================================================

\documentclass[letterpaper,11pt]{article}

\usepackage{latexsym}
\usepackage[empty]{fullpage}
\usepackage{titlesec}
\usepackage{marvosym}
\usepackage[usenames,dvipsnames]{color}
\usepackage{verbatim}
\usepackage{enumitem}
\usepackage[hidelinks]{hyperref}
\usepackage{fancyhdr}
\usepackage[english]{babel}
\usepackage{tabularx}
\usepackage{fontawesome5}
\usepackage{tikz}
\input{glyphtounicode}
\usepackage[margin=1.4cm]{geometry}

\pagestyle{fancy}
\fancyhf{} % clear all header and footer fields
\fancyfoot{}
\renewcommand{\headrulewidth}{0pt}
\renewcommand{\footrulewidth}{0pt}

% Adjust margins
\addtolength{\oddsidemargin}{-0.15in}
\addtolength{\textwidth}{0.3in}

\urlstyle{same}

\raggedbottom
\raggedright
\setlength{\tabcolsep}{0in}

% Sections formatting -- identical to resume.tex
\titleformat{\section}{
  \vspace{-4pt}\scshape\raggedright\large\bfseries
}{}{0em}{}[\color{black}\titlerule \vspace{-5pt}]

% Ensure that the generated pdf is machine readable/ATS parsable
\pdfgentounicode=1

%-------------------------
% Custom commands (reused from resume.tex)
\newcommand{\resumeItem}[1]{
  \item\small{
    {#1 \vspace{0pt}}
  }
}

\renewcommand\labelitemi{$\vcenter{\hbox{\tiny$\bullet$}}$}
\renewcommand\labelitemii{$\vcenter{\hbox{\tiny$\bullet$}}$}

\newcommand{\resumeItemListStart}{\begin{itemize}[leftmargin=0.18in]}
\newcommand{\resumeItemListEnd}{\end{itemize}\vspace{-5pt}}

% Tech stack tag -- small rounded chip, grayscale to match the
% monochrome resume aesthetic (no color usage beyond black/gray)
\newcommand{\techtag}[1]{%
  \tikz[baseline=(t.base)]{
    \node[draw=black, rounded corners=2pt, inner sep=3pt, fill=black!5]
      (t) {\small #1};
  }%
}

% Result callout box -- thin black border, mirrors the resume's
% \titlerule for visual consistency
\newcommand{\resultbox}[1]{%
  \noindent
  \fcolorbox{black}{black!4}{%
    \begin{minipage}{\dimexpr\linewidth-2\fboxsep-2\fboxrule}
      \small #1
    \end{minipage}%
  }%
}

\begin{document}

%----------HEADER----------
\begin{center}
    {\Large \scshape [PROJECT NAME]} \\[2mm]
    \normalsize \textbf{[YOUR ROLE]} $\bullet$ [DATE RANGE, e.g. Jan 2026 -- Mar 2026] \\[1.5mm]
    \footnotesize
    {\faGithub\ \underline{\href{[CODE URL]}{[code/demo link text]}}} ~~
    {\faBriefcase\ \underline{\href{[LIVE DEMO URL]}{[live demo link text]}}}
\end{center}
\vspace{-6pt}

%-----------PROBLEM-----------
\section{Problem}
\vspace{2pt}
\noindent\small
[2-3 sentence problem statement. Describe the pain point, who it affected,
and why solving it mattered. Keep this tight -- it sets up the Approach
and Result sections below.]

\vspace{6pt}

%-----------APPROACH-----------
\section{Approach}
    \resumeItemListStart
        \resumeItem{[Bullet describing a key design/technical decision and why you made it.]}
        \resumeItem{[Bullet describing an implementation detail, architecture choice, or tool integration.]}
        \resumeItem{[Bullet describing how you validated, tested, or iterated on the solution.]}
        \resumeItem{[Optional fourth bullet -- collaboration, constraints, or trade-offs handled.]}
    \resumeItemListEnd

%-----------RESULT-----------
\section{Result}
\vspace{4pt}
\resultbox{\textbf{[Headline outcome, quantified where possible -- e.g. ``Reduced page load time by \textbf{40\%} and onboarded \textbf{500+} users in the first month.'']} \\[2pt]
[Optional second sentence with a secondary metric or qualitative impact.]}
\vspace{4pt}

%-----------TECH STACK-----------
\section{Tech Stack}
\vspace{6pt}
\noindent
\techtag{[Language/Framework]} \hspace{4pt}
\techtag{[Language/Framework]} \hspace{4pt}
\techtag{[Tool/Platform]} \hspace{4pt}
\techtag{[Tool/Platform]} \hspace{4pt}
\techtag{[API/Service]} \hspace{4pt}
\techtag{[API/Service]}

\end{document}
